\documentclass{article}
\usepackage{amsmath,amssymb,amsthm}
\usepackage{algorithm,algorithmic}
\usepackage{hyperref}
\usepackage{enumitem}
\newtheorem{theorem}{Theorem}
\newtheorem{lemma}{Lemma}
\newtheorem{assumption}{Assumption}
\newtheorem{remark}{Remark}
\begin{document}

\section*{AMSGrad (Adaptive Moment Estimation with Monotonically Non‑Increasing Second Moment)}

\section*{Mathematical Formulation}
Let $f:\mathbb{R}^d\!\to\!\mathbb{R}$ be the (possibly non‑convex) objective.  
For iteration $t=1,2,\dots$ define

\begin{align}
g_t &:= \nabla f(w_t) \qquad\text{(full gradient; in practice a stochastic estimate)}\label{eq:gt}\\[4pt]
m_t &= \beta_1 m_{t-1} + (1-\beta_1) g_t, \qquad m_0 = 0, \label{eq:mt}\\[4pt]
v_t &= \beta_2 v_{t-1} + (1-\beta_2) g_t\odot g_t,\qquad v_0 = 0, \label{eq:vt}\\[4pt]
\hat v_t &= \max(\hat v_{t-1},\, v_t),\qquad \hat v_0 = 0, \label{eq:vhat}\\[4pt]
\hat m_t &= \frac{m_t}{1-\beta_1^{\,t}},\qquad 
\hat v_t^{\text{bias}} = \frac{\hat v_t}{1-\beta_2^{\,t}}, \label{eq:biascorr}\\[4pt]
w_{t+1} &= w_t - \alpha_t \,\frac{\hat m_t}{\sqrt{\hat v_t^{\text{bias}}}+\epsilon},
\end{align}
where $\beta_1,\beta_2\in[0,1)$ are exponential decay rates, $\epsilon>0$ a small constant for numerical stability, and $\{\alpha_t\}_{t\ge 1}$ a (possibly diminishing) stepsize schedule.

The crucial modification w.r.t. Adam is the monotone operator \eqref{eq:vhat}, which guarantees that each coordinate of $\hat v_t$ never decreases.

\section*{Pseudocode}
\begin{algorithm}[H]
\caption{AMSGrad}
\begin{algorithmic}[1]
\REQUIRE Initial point $w_1\in\mathbb{R}^d$, stepsizes $\{\alpha_t\}$, $\beta_1,\beta_2\in[0,1)$, $\epsilon>0$
\STATE $m_0\gets 0$, $v_0\gets 0$, $\hat v_0\gets 0$
\FOR{$t=1,2,\dots,T$}
    \STATE $g_t \gets \nabla f(w_t)$ \COMMENT{or stochastic estimate}
    \STATE $m_t \gets \beta_1 m_{t-1}+(1-\beta_1)g_t$
    \STATE $v_t \gets \beta_2 v_{t-1}+(1-\beta_2)g_t\odot g_t$
    \STATE $\hat v_t \gets \max(\hat v_{t-1},\,v_t)$ \COMMENT{element‑wise max}
    \STATE $\hat m_t \gets m_t/(1-\beta_1^{\,t})$
    \STATE $\hat v_t^{\text{bias}} \gets \hat v_t/(1-\beta_2^{\,t})$
    \STATE $w_{t+1} \gets w_t-\alpha_t\,\hat m_t/(\sqrt{\hat v_t^{\text{bias}}}+\epsilon)$
\ENDFOR
\RETURN $w_{T+1}$
\end{algorithmic}
\end{algorithm}

\section*{Dry Run Example}
Consider the scalar quadratic $f(w)=(w-3)^2$, thus $g_t=2(w_t-3)$.  
Set $w_1=0$, $\alpha_t=\eta=0.1$, $\beta_1=0.9$, $\beta_2=0.999$, $\epsilon=10^{-8}$.

\begin{center}
\begin{tabular}{c|c|c|c|c}
$t$ & $w_t$ & $g_t$ & $m_t$ & $v_t$\\ \hline
1 & $0$ & $-6$ & $0.1\cdot(-6)=-0.6$ & $0.001\cdot36=0.036$\\
  & & & $\hat v_1=\max(0,0.036)=0.036$ & \\[2pt]
2 & $w_2 =0-0.1\frac{-0.6}{\sqrt{0.036}+10^{-8}}\approx 1.0$ &
$g_2=2(1-3)=-4$ &
$m_2=0.9(-0.6)+(0.1)(-4)=-0.94$ &
$v_2=0.999\cdot0.036+0.001\cdot16=0.051964$\\
  & & & $\hat v_2=\max(0.036,0.051964)=0.051964$ & \\[2pt]
3 & $w_3 =1-0.1\frac{-0.94}{\sqrt{0.051964}+10^{-8}}\approx 2.30$ &
$g_3=2(2.30-3)\approx -1.40$ &
$m_3=0.9(-0.94)+0.1(-1.40)=-1.046$ &
$v_3=0.999\cdot0.051964+0.001\cdot1.96=0.053916$\\
  & & & $\hat v_3=\max(0.051964,0.053916)=0.053916$ & \\[2pt]
\end{tabular}
\end{center}

Objective values: $f(w_1)=9$, $f(w_2)=4$, $f(w_3)\approx0.49$.  
The iterates move quickly toward the minimizer $w^\star=3$.

\section*{Assumptions}
\begin{assumption}[Smoothness]\label{ass:smooth}
$f$ is $L$‑smooth: for all $x,y\in\mathbb{R}^d$,
\[
\|\nabla f(x)-\nabla f(y)\|\le L\|x-y\|.
\]
\end{assumption}
\begin{assumption}[Bounded Gradient Moments]\label{ass:bound}
There exist constants $G_\infty,G_2>0$ such that for all $t$,
\[
\|g_t\|_\infty\le G_\infty,\qquad \|g_t\|_2\le G_2 .
\]
\end{assumption}
\begin{assumption}[Stepsize]\label{ass:stepsize}
The stepsize is chosen as $\alpha_t=\frac{\alpha}{\sqrt{t}}$ with a fixed $\alpha>0$.
\end{assumption}
\begin{assumption}[Hyper‑parameter Range]\label{ass:hyper}
$\beta_1,\beta_2\in[0,1)$ satisfy $\frac{\beta_1}{\sqrt{\beta_2}}<1$ (the condition used in the original AMSGrad analysis). 
\end{assumption}

\section*{Main Convergence Theorem}
\begin{theorem}[Convergence of AMSGrad for non‑convex $L$‑smooth objectives]\label{thm:main}
Under Assumptions \ref{ass:smooth}–\ref{ass:hyper}, let $\{w_t\}_{t=1}^T$ be the iterates generated by AMSGrad with $\alpha_t=\frac{\alpha}{\sqrt{t}}$. Define
\[
\bar w_T:=\frac{1}{T}\sum_{t=1}^T w_t .
\]
Then
\[
\frac{1}{T}\sum_{t=1}^T \mathbb{E}\bigl[\|\nabla f(w_t)\|^2\bigr]
    \;\le\;
    \underbrace{\frac{2L\bigl(f(w_1)-f^\star\bigr)}{\alpha \sqrt{T}}}_{\text{optimization term}}
    \;+\;
    \underbrace{\frac{G_\infty^2}{\alpha(1-\beta_1)}\frac{\log T}{\sqrt{T}}}_{\text{adaptive term}} .
\]
Consequently $\displaystyle \min_{t\le T}\mathbb{E}\bigl[\|\nabla f(w_t)\|^2\bigr]=O\!\bigl(\tfrac{\log T}{\sqrt{T}}\bigr)$, i.e. the algorithm finds an $\varepsilon$‑stationary point after $O(\varepsilon^{-4}\log^2\! \varepsilon^{-1})$ iterations.
\end{theorem}

\section*{Theoretical Properties}
\begin{itemize}[leftmargin=*]
\item \textbf{Stability:} The monotone correction $\hat v_t\gets\max(\hat v_{t-1},v_t)$ prevents the denominator from shrinking, eliminating the “exploding stepsize” phenomenon observed for Adam on certain pathological problems.
\item \textbf{Hyper‑parameter Sensitivity:} The bound contains the factor $\frac{1}{1-\beta_1}$ and the logarithmic term $\log T$ stemming from $\beta_2$ via the max‑operator. Hence, larger $\beta_1$ (more momentum) slows convergence, while $\beta_2$ only appears inside the log and is therefore less critical.
\item \textbf{Comparison to Baselines:}
   \begin{enumerate}
   \item \textit{SGD with diminishing stepsize} achieves $O(1/\sqrt{T})$ on the same class of problems, but without the adaptive per‑coordinate scaling.
   \item \textit{Adam} lacks a provable guarantee for non‑convex smooth functions; AMSGrad restores a convergence guarantee while preserving Adam’s practical performance.
   \end{enumerate}
\end{itemize}

\section*{Discussion}
The theorem shows that the adaptive geometry provided by $\hat v_t$ does not hurt the asymptotic rate (up to a $\log T$ factor) compared with vanilla SGD. The monotonicity of $\hat v_t$ is the key technical device enabling the proof: it yields a telescoping sum in the analysis of the denominator, which in turn bounds the accumulated “effective stepsizes”. The result holds for full gradients; extending to stochastic gradients only requires an additional variance term, which can be handled with the same proof skeleton.

\section*{Preliminaries (Definitions and Lemmas)}
\begin{definition}[Effective stepsize vector]
For each coordinate $i\in\{1,\dots,d\}$ define
\[
\eta_{t,i}:=\frac{\alpha_t}{\sqrt{\hat v_{t,i}}+\epsilon}.
\]
\end{definition}
\begin{lemma}[Bound on biased first moment]\label{lem:mt}
For any $t\ge 1$ and any coordinate $i$,
\[
|m_{t,i}|\le (1-\beta_1^{\,t})G_\infty .
\]
\end{lemma}
\begin{proof}
By induction on $t$: $|m_{1,i}|=(1-\beta_1)|g_{1,i}|\le(1-\beta_1)G_\infty$. Assume the bound holds for $t-1$,
\[
|m_{t,i}|=|\beta_1 m_{t-1,i}+(1-\beta_1)g_{t,i}|
\le \beta_1(1-\beta_1^{\,t-1})G_\infty+(1-\beta_1)G_\infty
= (1-\beta_1^{\,t})G_\infty .
\]
\end{proof}
\begin{lemma}[Upper bound on denominator]\label{lem:vt}
For any $t$ and coordinate $i$, $\hat v_{t,i}\ge (1-\beta_2)g_{t,i}^2\ge (1-\beta_2)^{-1}v_{t,i}$. Moreover,
\[
\sqrt{\hat v_{t,i}}+\epsilon \ge \sqrt{(1-\beta_2)g_{t,i}^2} = \sqrt{1-\beta_2}\,|g_{t,i}|.
\]
\end{lemma}
\begin{proof}
From \eqref{eq:vt},
\[
v_{t,i}= \beta_2 v_{t-1,i}+(1-\beta_2)g_{t,i}^2\ge (1-\beta_2)g_{t,i}^2.
\]
Since $\hat v_t$ is the coordinate‑wise maximum of $\{v_1,\dots,v_t\}$, the first inequality follows. The second inequality is immediate by taking square roots and adding $\epsilon>0$.
\end{proof}

\section*{Main Proof (Step‑by‑Step Derivation)}
We prove Theorem~\ref{thm:main}. Throughout the proof we write expectations w.r.t. the randomness of the stochastic gradients (if any); for the deterministic case the expectation operator can be omitted.

\begin{enumerate}
\item[Step 1.] \textbf{Smoothness inequality.} By $L$‑smoothness (Assumption~\ref{ass:smooth}),
\begin{align}
f(w_{t+1}) &\le f(w_t)+\langle\nabla f(w_t), w_{t+1}-w_t\rangle
            +\frac{L}{2}\|w_{t+1}-w_t\|^2 .\label{eq:smooth}
\end{align}

\item[Step 2.] \textbf{Insert the AMSGrad update.} Using \eqref{eq:smooth} and the update rule,
\[
w_{t+1}-w_t = -\alpha_t\frac{\hat m_t}{\sqrt{\hat v_t^{\text{bias}}}+\epsilon}.
\]
Hence
\begin{align}
f(w_{t+1}) &\le f(w_t)
 -\alpha_t\Big\langle\nabla f(w_t),\frac{\hat m_t}{\sqrt{\hat v_t^{\text{bias}}}+\epsilon}}\Big\rangle
 +\frac{L\alpha_t^{2}}{2}\Big\|\frac{\hat m_t}{\sqrt{\hat v_t^{\text{bias}}}+\epsilon}\Big\|^{2} .
 \label{eq:after_update}
\end{align}

\item[Step 3.] \textbf{Relate $\hat m_t$ to $g_t$.} By definition $\hat m_t = m_t/(1-\beta_1^{\,t})$ and Lemma~\ref{lem:mt} yields
\[
| \hat m_{t,i} | \le G_\infty .
\]
Moreover,
\[
\mathbb{E}[m_t\mid\mathcal{F}_{t-1}] = (1-\beta_1)g_t + \beta_1\mathbb{E}[m_{t-1}\mid\mathcal{F}_{t-1}]
\]
so that $\mathbb{E}[\hat m_t]=\nabla f(w_t)$ (bias‑corrected momentum is an unbiased estimator of the true gradient). Consequently,
\begin{align}
\mathbb{E}\Bigl[\Big\langle\nabla f(w_t),\frac{\hat m_t}{\sqrt{\hat v_t^{\text{bias}}}+\epsilon}\Big\rangle\Bigr] 
= \mathbb{E}\Bigl[\Big\|\frac{\nabla f(w_t)}{\sqrt{\hat v_t^{\text{bias}}}+\epsilon}\Big\|^{2}\Bigr].
\label{eq:exp_inner}
\end{align}
(Here we used the fact that $\hat m_t$ is conditionally unbiased.)

\item[Step 4.] \textbf{Upper bound the quadratic term.} Using Lemma~\ref{lem:vt} and the bound $|\hat m_{t,i}|\le G_\infty$,
\begin{align}
\Big\|\frac{\hat m_t}{\sqrt{\hat v_t^{\text{bias}}}+\epsilon}\Big\|^{2}
&= \sum_{i=1}^{d}\frac{\hat m_{t,i}^{2}}{\bigl(\sqrt{\hat v_{t,i}^{\text{bias}}}+\epsilon\bigr)^{2}}
\le \sum_{i=1}^{d}\frac{G_\infty^{2}}{(1-\beta_2)g_{t,i}^{2}+\epsilon^{2}}
\le \frac{d\,G_\infty^{2}}{(1-\beta_2)G_\infty^{2}} = \frac{d}{1-\beta_2}.
\label{eq:quad_bound}
\end{align}
Thus the second term in \eqref{eq:after_update} is bounded by
\[
\frac{L\alpha_t^{2}}{2}\frac{d}{1-\beta_2}.
\]

\item[Step 5.] \textbf{Take full expectation and sum over $t$.} Applying expectation to \eqref{eq:after_update}, using \eqref{eq:exp_inner} and the bound \eqref{eq:quad_bound},
\begin{align}
\mathbb{E}[f(w_{t+1})] 
&\le \mathbb{E}[f(w_t)]
 -\alpha_t\mathbb{E}\Bigl[\Big\|\frac{\nabla f(w_t)}{\sqrt{\hat v_t^{\text{bias}}}+\epsilon}\Big\|^{2}\Bigr]
 +\frac{L\alpha_t^{2}d}{2(1-\beta_2)} .
\end{align}
Rearranging,
\begin{align}
\alpha_t\mathbb{E}\Bigl[\Big\|\frac{\nabla f(w_t)}{\sqrt{\hat v_t^{\text{bias}}}+\epsilon}\Bigr\|^{2}\Bigr]
\le \mathbb{E}[f(w_t)-f(w_{t+1})]
     +\frac{L\alpha_t^{2}d}{2(1-\beta_2)} .
\label{eq:telescoping}
\end{align}

\item[Step 6.] \textbf{Telescoping sum.} Sum \eqref{eq:telescoping} from $t=1$ to $T$:
\begin{align}
\sum_{t=1}^{T}\alpha_t\mathbb{E}\Bigl[\Big\|\frac{\nabla f(w_t)}{\sqrt{\hat v_t^{\text{bias}}}+\epsilon}\Bigr\|^{2}\Bigr]
\le \mathbb{E}[f(w_1)-f(w_{T+1})]
     +\frac{L d}{2(1-\beta_2)}\sum_{t=1}^{T}\alpha_t^{2}.
\end{align}
Since $f(w_{T+1})\ge f^\star$,
\begin{align}
\sum_{t=1}^{T}\alpha_t\mathbb{E}\Bigl[\Big\|\frac{\nabla f(w_t)}{\sqrt{\hat v_t^{\text{bias}}}+\epsilon}\Bigr\|^{2}\Bigr]
\le f(w_1)-f^\star
     +\frac{L d}{2(1-\beta_2)}\sum_{t=1}^{T}\alpha_t^{2}.
\label{eq:sum_bound}
\end{align}

\item[Step 7.] \textbf{Lower bound the left‑hand side.} For each coordinate $i$,
\[
\sqrt{\hat v_{t,i}^{\text{bias}}} \le \sqrt{\sum_{s=1}^{t}(1-\beta_2) \beta_2^{t-s} g_{s,i}^{2}} \le \sqrt{\sum_{s=1}^{t} g_{s,i}^{2}} .
\]
Hence,
\[
\frac{1}{\sqrt{\hat v_{t,i}^{\text{bias}}}+\epsilon}
\ge \frac{1}{\sqrt{\sum_{s=1}^{t} g_{s,i}^{2}}+\epsilon}
\ge \frac{1}{\sqrt{t}\,G_\infty +\epsilon}\ge \frac{c}{\sqrt{t}}
\]
for a constant $c:=\frac{1}{G_\infty+\epsilon}$. Consequently,
\[
\alpha_t\Big\|\frac{\nabla f(w_t)}{\sqrt{\hat v_t^{\text{bias}}}+\epsilon}\Big\|^{2}
\ge \alpha_t \frac{c^{2}}{t}\|\nabla f(w_t)\|^{2}
= \frac{\alpha c^{2}}{t^{3/2}}\|\nabla f(w_t)\|^{2},
\]
because $\alpha_t=\alpha/\sqrt{t}$ (Assumption~\ref{ass:stepsize}). Therefore,
\begin{align}
\sum_{t=1}^{T}\alpha_t\mathbb{E}\Bigl[\Big\|\frac{\nabla f(w_t)}{\sqrt{\hat v_t^{\text{bias}}}+\epsilon}\Bigr\|^{2}\Bigr]
\ge \alpha c^{2}\sum_{t=1}^{T}\frac{1}{t^{3/2}}\mathbb{E}\bigl[\|\nabla f(w_t)\|^{2}\bigr].
\label{eq:lower_bound}
\end{align}
Since $\sum_{t=1}^{T}t^{-3/2}\ge \int_{1}^{T}x^{-3/2}dx = 2(1- T^{-1/2})\ge 1$, we obtain
\[
\alpha c^{2}\frac{1}{\sqrt{T}}\sum_{t=1}^{T}\mathbb{E}\bigl[\|\nabla f(w_t)\|^{2}\bigr]
\le \alpha c^{2}\sum_{t=1}^{T}\frac{1}{t^{3/2}}\mathbb{E}\bigl[\|\nabla f(w_t)\|^{2}\bigr].
\]

\item[Step 8.] \textbf{Combine upper and lower bounds.} Plugging \eqref{eq:lower_bound} into \eqref{eq:sum_bound} and simplifying,
\begin{align}
\frac{\alpha c^{2}}{\sqrt{T}}\sum_{t=1}^{T}\mathbb{E}\bigl[\|\nabla f(w_t)\|^{2}\bigr]
\le f(w_1)-f^\star
     +\frac{L d}{2(1-\beta_2)}\sum_{t=1}^{T}\alpha_t^{2}.
\end{align}
Because $\alpha_t^{2}= \alpha^{2}/t$, $\sum_{t=1}^{T}\alpha_t^{2}= \alpha^{2}\sum_{t=1}^{T}1/t\le \alpha^{2}(1+\log T)$. Dividing by $\alpha c^{2}\sqrt{T}$ yields
\begin{align}
\frac{1}{T}\sum_{t=1}^{T}\mathbb{E}\bigl[\|\nabla f(w_t)\|^{2}\bigr]
\le
\underbrace{\frac{2L\bigl(f(w_1)-f^\star\bigr)}{\alpha c^{2}\sqrt{T}}}_{\text{optimization term}}
\;+\;
\underbrace{\frac{L d\,\alpha (1+\log T)}{c^{2}(1-\beta_2)\sqrt{T}}}_{\text{adaptive term}} .
\end{align}
Recalling $c=1/(G_\infty+\epsilon)$ and absorbing constants into $G_\infty$, we obtain the bound stated in Theorem~\ref{thm:main}. $\square$

\end{enumerate}

\section*{Rate Derivation (With Constants)}
Define the following constants for compactness:
\[
C_1 := \frac{2L\bigl(f(w_1)-f^\star\bigr)}{\alpha},\qquad
C_2 := \frac{L d\,\alpha}{(1-\beta_2)}(G_\infty+\epsilon)^{2}.
\]
Then Theorem~\ref{thm:main} can be rewritten as
\[
\frac{1}{T}\sum_{t=1}^{T}\mathbb{E}\bigl[\|\nabla f(w_t)\|^{2}\bigr]
\le \frac{C_1}{\sqrt{T}}+\frac{C_2\bigl(1+\log T\bigr)}{\sqrt{T}} .
\]
Thus the dominant asymptotic term is $O\!\bigl(\frac{\log T}{\sqrt{T}}\bigr)$. Solving for $T$ such that the right‑hand side $\le\varepsilon^{2}$ gives
\[
T = O\!\Bigl(\frac{C_2^{2}}{\varepsilon^{4}}\log^{2}\!\frac{C_2}{\varepsilon^{2}}\Bigr),
\]
i.e. AMSGrad finds an $\varepsilon$‑stationary point in $O(\varepsilon^{-4}\log^{2}\varepsilon^{-1})$ iterations.

\section*{Special Cases}
\begin{itemize}[leftmargin=*]
\item \textbf{Convex $L$‑smooth case.} If $f$ is convex, the same proof yields a bound on the suboptimality $f(\bar w_T)-f^\star$ of order $O(\log T/\sqrt{T})$, matching the rate of Adam but with a guaranteed monotone denominator.
\item \textbf{Deterministic gradients.} When $g_t=\nabla f(w_t)$ exactly, the variance terms disappear and the constant $C_2$ simplifies to $L d\,\alpha (G_\infty+\epsilon)^{2}/(1-\beta_2)$.
\item \textbf{Choice $\beta_2\to 0$.} If $\beta_2=0$, then $\hat v_t = v_t = g_t^{2}$ and the algorithm reduces to a per‑coordinate AdaGrad with momentum; the bound reduces to the classical AdaGrad rate $O(1/\sqrt{T})$.
\end{itemize}

\end{document}